The cap at $b$ reflects that a renovated seller keeps her hiding cost ($c_1=0$): the rent that renovation destroys cannot exceed the rating premium. If the subsidy required her to display the new rating, the cap would be absent and the bound would be $\hat\sigma_0\min\{\lambda,\hat R\}$. At the estimates the bound is \euro{}\attSavingUbTwoFiftyApt{} billion per year in apartments, or \euro{}\attSavingUbDTwoFiftyApt{} billion at the light benchmark and \euro{}\attSavingUbDFiveHundredApt{} billion at the medium one if the renovated seller must display. In houses it is \euro{}\attSavingUbTwoFiftyHouse{} billion, because $\hat\psi$ is close to one and the price of silence barely falls (\attSalienceHouse{} log points). Together the bound is \attSavingUbTotTwoFiftyLo{} to \attSavingUbTotTwoFiftyHi{} percent of MaPrimeR\'enov's annual budget, and the largest apartment market accounts for \attSavingTopShareApt{} percent of it. Using the observed wedge in place of the model's price of silence gives \euro{}\attSavingUbObsTwoFiftyApt{} billion for apartments. The hiding rate behind these figures is the maintained one, \attSigmaObsApt{} percent, which assigns every informed hider to low quality and so also counts hiders whose certificate is A to D; with the rate computed from E to G certificates alone, \attSigmaRegMeanApt{} percent, the bound is \euro{}\attSavingUbRegTwoFiftyApt{} billion.
